\section{Intermediate-Guided Reasoning：把推理过程变成可操作对象}

第二个单元从“能力何时出现”转向“怎样组织一次推理”。讲者明确说 Intermediate-Guided Reasoning 不是标准术语，而是对一族方法的概括：让模型生成自然语言步骤、搜索多个 thought states、递归分解问题，或把中间表示交给程序解释器。关键区别在于中间状态能否被验证、回退和执行。

\subsection{一个课堂 umbrella term}

标题页标记转场，真正定义来自口头解释：这些方法把最终答案前的 intermediate information 显式化，使复杂任务可以被拆分、检查或交给外部模块。本节先说明这个 umbrella term 的课堂语境，再把它拆成文本链、搜索树、程序和计算图四类表示。

\lecturefigure{slide-13-intermediate-reasoning.jpg}{Intermediate-Guided Reasoning 课堂单元}{Stanford 课堂视频 00:07:42}

\begin{knowledgebox}{读图：不要把课堂短语误写成固定 taxonomy}
本讲把 CoT、ToT、Socratic decomposition、PAL 和 PoT 放在同一伞下，是为了比较 intermediate representation；学术文献并没有统一采用这个总称。讲义保留这种教学组织，同时继续使用各论文原名。
\end{knowledgebox}

\teachervoice{讲者主动说“I don't think this is actually a term”。这是必须保留的 teacher voice：术语来源是课堂组织，不应包装成公认领域名称。}

\subsection{Chain-of-Thought 的基本机制}

CoT 让模型在最终答案前生成一系列自然语言 reasoning steps。对多步 arithmetic、symbolic 和 commonsense tasks，它可以把一次困难映射拆成多个较短条件生成；但生成更多 token 也增加传播错误和制造貌似合理解释的机会。

\lecturefigure{slide-14-cot-overview.jpg}{Chain-of-Thought 以中间文本分解多步问题}{Stanford 课堂视频 00:08:03}

\begin{knowledgebox}{读图：CoT 改变 inference trace 而非模型权重}
Few-shot CoT 通常在 prompt 中展示“问题→步骤→答案”样例，模型随后模仿格式。它可能调动训练中已有的算法模式，但不能保证每一步对应真实内部因果过程。能力提升与解释忠实性是两个不同验收项。
\end{knowledgebox}

给定问题 $x$，中间轨迹 $z=(z_1,\ldots,z_K)$，答案 $y$，CoT 生成可写为
\[
p_\theta(y,z\mid x)=p_\theta(z\mid x)p_\theta(y\mid x,z).
\]
$z$ 表示可见 reasoning trace，$K$ 表示步骤数，$y$ 表示最终答案。提高 $p(y,z\mid x)$ 不等于 $z$ 忠实解释模型内部计算，只说明该文本轨迹与答案共同高概率。

\subsection{Standard prompting 与 CoT prompting}

例子页用网球问题对比直接输出和逐步计算。读者应观察的不只是答案是否正确，还要检查步骤是否保持单位、是否遗漏数量、最后答案是否真正由前面步骤推出。本节从同一道题的两种 prompt 入手，说明可见轨迹怎样改变诊断能力，却不自动构成 proof。

\lecturefigure{slide-15-cot-examples.jpg}{Standard prompting 与 Chain-of-Thought prompting 对比}{Stanford 课堂视频 00:09:18}

\begin{knowledgebox}{读图：步骤提供了调试表面}
Direct answer 错误时很难判断模型在哪里偏离；CoT 至少暴露了 arithmetic、symbol mapping、missing step 或 semantic misunderstanding。可见性提高了 diagnosis，却不自动提高 truthfulness，因为模型也可能先得到答案再编写理由。
\end{knowledgebox}

\subsection{错误分析比平均准确率更有用}

Improving CoT slide 把错误拆成 calculator error、symbol mapping、missing step 与 incoherent reasoning 等类别。讲者强调，即使大模型获得明显 performance gain，non-negligible errors 仍然存在，因此需要工具、verifier 与更好的训练。

\lecturefigure{slide-16-cot-improvements.jpg}{Chain-of-Thought 的错误来源与改进方向}{Stanford 课堂视频 00:09:42}

\begin{knowledgebox}{读图：错误类别对应不同修复层}
Calculator error 可交给 interpreter；symbol mapping 需要更稳定的变量绑定；missing step 可用 decomposition 或 search；incoherent CoT 需要 verifier、consistency 或过程监督。把所有错误都归因于“模型不够大”会错过低成本修复。
\end{knowledgebox}

\begin{warningbox}{Reasoning trace 不是 proof}
一条流畅 CoT 可能包含无关步骤、隐含跳跃或后验合理化。高风险任务要验证 executable constraints、引用证据或环境结果，而不是用自然语言长度代替正确性。
\end{warningbox}

\subsection{怎样让小模型也使用 CoT}

课堂指出，早期 CoT 对约百亿到百亿以上规模更有效，小模型常出现 missing step 与 semantic understanding 错误。研究方向包括用大模型生成 traces 做 fine-tuning、distillation、更精简的 reasoning format 和专门 verifier。

\lecturefigure{slide-17-cot-smaller-models.jpg}{面向小模型的 Chain-of-Thought 研究方向}{Stanford 课堂视频 00:10:18}

\begin{knowledgebox}{读图：小模型问题可能是表示容量也可能是监督格式}
若 trace 太长或包含噪声，小模型会先学到表面模板；若任务需要无法表示的 latent state，单纯蒸馏也不足。应分别评估 final answer、step validity、trace length 与 compute，避免把“大模型 trace”视为天然最优教材。
\end{knowledgebox}

\subsection{从特定领域走向更一般的 reasoning format}

Generalization slide 质疑 CoT 的固定格式与领域依赖：它在多步数学问题上突出，但并非所有任务都适合线性自然语言步骤。更一般的方法需要允许不同 state representation、branching、external calls 和 stopping rules。

\lecturefigure{slide-18-cot-generalization.jpg}{Chain-of-Thought 的格式刚性与领域泛化问题}{Stanford 课堂视频 00:11:03}

\begin{knowledgebox}{读图：推理表示应匹配任务结构}
线性文本适合顺序分解；规划任务可能需要树或图；精确算术适合程序；感知任务可能需要视觉 state。Generalization 不是强迫所有任务输出同一种长解释，而是选择能被验证和执行的 intermediate representation。
\end{knowledgebox}

\subsection{Tree of Thoughts：从一条轨迹到状态搜索}

Tree of Thoughts 让模型生成多个候选 thought，评估下一步，必要时 look ahead 或 backtrack。它把推理从一次 autoregressive sample 改造成 search procedure，因此获得更高 flexibility，也付出更多 inference calls 与 evaluator dependence。

\lecturefigure{slide-19-tree-of-thoughts.jpg}{Tree of Thoughts 的分支评估与回溯}{Stanford 课堂视频 00:12:03}

\begin{knowledgebox}{读图：节点是中间状态，边是推理动作}
同一 input 产生多个 candidate states，evaluator 给分后保留部分分支；失败分支可回退。图中最关键的不是树形外观，而是 generation、evaluation、selection 三个角色被分开，允许显式控制搜索。
\end{knowledgebox}

若每层扩展 $b$ 个候选、深度为 $d$，完整搜索最多访问
\[
1+b+b^2+\cdots+b^d=\frac{b^{d+1}-1}{b-1}
\]
个节点。$b$ 是 branching factor，$d$ 是 reasoning depth。ToT 必须用 beam、heuristic 或 pruning 控制指数开销；evaluator 错误也会提前剪掉正确分支。

\subsection{Socratic Questioning：递归分解与回收答案}

Socratic Questioning 用大模型提出 subproblems，递归求解后再向上组合。它与 ToT 都允许 backtracking，但教学重点更接近 divide-and-conquer：先确认子问题是否可解，再把局部答案汇总为原问题解。

\lecturefigure{slide-20-socratic-questioning.jpg}{Socratic Questioning 的递归子问题分解}{Stanford 课堂视频 00:12:33}

\begin{knowledgebox}{读图：decomposition quality 决定上限}
若子问题不完备、彼此重叠或丢失约束，即使每个局部答案正确也无法得到全局答案。系统应检查 coverage、dependency 和 recomposition，而不是只检查每个 leaf 是否看似合理。
\end{knowledgebox}

\begin{lstlisting}[caption={带验证与回溯的结构化推理流程}]
def solve(problem, depth):
    if depth == 0 or is_atomic(problem):
        return answer_and_verify(problem)
    subproblems = propose_decomposition(problem)
    partial = [solve(item, depth - 1) for item in subproblems]
    candidate = compose(problem, partial)
    if verify(candidate, problem):
        return candidate
    return revise_decomposition(problem, partial)
\end{lstlisting}

\subsection{PAL：语言模型规划，解释器计算}

Program-Aided Language Models 让模型生成程序，把精确 arithmetic 和 symbolic execution 交给 Python 等 interpreter。LLM 负责把自然语言转成可执行结构，计算机负责确定性运行；这比要求模型在 token space 内模拟全部算术更可靠。

\lecturefigure{slide-21-pal.jpg}{PAL 用程序作为中间表示并调用解释器}{Stanford 课堂视频 00:13:09}

\begin{knowledgebox}{读图：把神经推理与符号执行分工}
Prompt 中展示问题、program 与 answer；模型输出 code，interpreter 运行得到结果。优势是 arithmetic 可验证、trace 可执行；风险是 code generation 可能误解语义、调用危险 API 或在 sandbox 外产生副作用。
\end{knowledgebox}

\subsection{Program of Thoughts：更一般的程序化 scratchpad}

Program of Thoughts 同样让模型生成 program，但强调把 computation 与 language reasoning 交错组织。对复杂数学题，程序保存变量、循环和函数关系，避免长自然语言链中的算术漂移。本节承接 PAL 的 interpreter 分工，进一步关注程序化 scratchpad 如何保存长程状态与组合结构。

\lecturefigure{slide-22-program-of-thoughts.jpg}{Program of Thoughts 的程序生成与执行流程}{Stanford 课堂视频 00:13:51}

\begin{knowledgebox}{读图：程序正确不等于问题理解正确}
Interpreter 只保证给定程序按语义执行；如果模型把问题映射成错误公式，程序会稳定地产生错误答案。因此 pipeline 至少需要 syntax check、unit test、constraint check 与 result sanity check。
\end{knowledgebox}

可把 program-guided reasoning 写成
\[
P\sim p_\theta(P\mid x),\qquad y=\mathcal{I}(P),
\]
其中 $P$ 是模型生成程序，$\mathcal{I}$ 是受控解释器，$y$ 是执行结果。可靠性取决于 semantic parsing $p_\theta$ 与 execution sandbox $\mathcal{I}$ 两部分，而不是只取决于 final answer accuracy。

\subsection{Faith and Fate：把 compositionality 写成计算图}

课堂最后展示 computation graph 工作，强调结构化中间表示可以把复杂问题拆成可检查节点。图中的 graph 不一定优于所有文本方法，但它提供 dependency、reuse 和 local verification，适合组合性任务。

\lecturefigure{slide-23-faith-and-fate.jpg}{Faith and Fate 工作中的 compositional computation graph}{Stanford 课堂视频 00:14:33}

\begin{knowledgebox}{读图：图结构让依赖显式化}
节点表示中间子计算，边表示依赖；若最终答案错，可以沿 graph 定位错误节点。与线性 CoT 相比，graph 更容易复用共享子结果，也更适合并行执行；代价是 graph construction 与 node semantics 本身需要学习或规则。
\end{knowledgebox}

\subsection{本章小结}

Intermediate-guided methods 的共同点是让中间状态可见、可搜索或可执行。CoT 提供线性文本，ToT/Socratic 提供 search/decomposition，PAL/PoT 提供程序执行，computation graph 提供依赖结构；可靠性取决于 verifier、interpreter 和反馈，而不是轨迹长度。

